\section{Encoding-Regime Separation}\label{sec:dichotomy}

The hardness results of Section~\ref{sec:hardness} apply to worst-case instances under the succinct encoding. This section states an encoding-regime separation: an explicit-state upper bound versus a succinct-encoding worst-case lower bound, and an intermediate query-access obstruction family. The complexity class is a property of the encoding measured against a fixed decision question (Remark~\ref{rem:question-vs-problem}), not a property of the question itself: the same sufficiency question admits polynomial-time algorithms under one encoding and worst-case exponential cost under another.

\paragraph{Model note.} Theorem~\ref{thm:dichotomy} has Part~1 in [E] (time polynomial in $|S|$) and Part~2 in [S+ETH] (time exponential in $n$). Proposition~\ref{prop:query-regime-obstruction} is [Q\_fin] (finite-state Opt-oracle core), while Propositions~\ref{prop:query-value-entry-lb} and~\ref{prop:query-state-batch-lb} are [Q\_bool] interface refinements (value-entry and state-batch). Relative to Definition~\ref{def:regime-simulation}, Propositions~\ref{prop:query-subproblem-transfer} and~\ref{prop:oracle-lattice-transfer} are explicit simulation-transfer instances. The [E] vs [S+ETH] split in Theorem~\ref{thm:dichotomy} is a same-question encoding/cost-model separation, not a bidirectional simulation claim between those regimes. These regimes are defined in Section~\ref{sec:encoding} and are not directly comparable as functions of one numeric input-length parameter.

\begin{theorem}[Explicit--Succinct Regime Separation]\label{thm:dichotomy}
Let $\mathcal{D} = (A, X_1, \ldots, X_n, U)$ be a decision problem with $|S| = N$ states. Let $k^*$ be the size of the minimal sufficient set.

\begin{enumerate}
\item \textbf{[E] Explicit-state upper bound:} Under the explicit-state encoding, SUFFICIENCY-CHECK is solvable in time polynomial in $N$ (e.g. $O(N^2|A|)$).

\item \textbf{[S+ETH] Succinct lower bound (worst case):} Assuming ETH, there exists a family of succinctly encoded instances with $n$ coordinates and minimal sufficient set size $k^* = n$ such that SUFFICIENCY-CHECK requires time $2^{\Omega(n)}$.
\end{enumerate}
\leanmeta{\LH{CC13},
\LH{CC15},
\LH{CC16}.}
\end{theorem}

\begin{proof}
\textbf{Part 1 (Explicit-state upper bound):} Under the explicit-state encoding, SUFFICIENCY-CHECK is decidable in time polynomial in $N$ by direct enumeration: compute $\Opt(s)$ for all $s\in S$ and then check all pairs $(s,s')$ with $s_I=s'_I$.

Equivalently, for any fixed $I$, the projection map $s\mapsto s_I$ has image size $|S_I|\le |S|=N$, so any algorithm that iterates over projection classes (or over all state pairs) runs in polynomial time in $N$. Thus, in particular, when $k^*=O(\log N)$, SUFFICIENCY-CHECK is solvable in polynomial time under the explicit-state encoding.

\paragraph{Remark (bounded coordinate domains).} In the general model $S=\prod_i X_i$, for a fixed $I$ one always has $|S_I|\le \prod_{i\in I}|X_i|$ (and $|S_I|\le N$). If the coordinate domains are uniformly bounded, $|X_i|\le d$ for all $i$, then $|S_I|\le d^{|I|}$.

\textbf{Part 2 (Succinct ETH lower bound, worst case):} A strengthened version of the TAUTOLOGY reduction used in Theorem~\ref{thm:sufficiency-conp} produces a family of instances in which the minimal sufficient set has size $k^* = n$: given a Boolean formula $\varphi$ over $n$ variables, we construct a decision problem with $n$ coordinates such that if $\varphi$ is not a tautology then \emph{every} coordinate is decision-relevant (thus $k^*=n$). This strengthening is mechanized in Lean (see Appendix~\ref{app:lean}). Under the Exponential Time Hypothesis (ETH) \cite{impagliazzo2001complexity,arora2009computational}, TAUTOLOGY requires time $2^{\Omega(n)}$ in the succinct encoding, so SUFFICIENCY-CHECK inherits a $2^{\Omega(n)}$ worst-case lower bound via this reduction.
\end{proof}

\begin{corollary}[Regime Separation (by Encoding)]
There is a clean separation between explicit-state tractability and succinct worst-case hardness (with respect to the encodings in Section~\ref{sec:encoding}):
\begin{itemize}
\item Under the explicit-state encoding, SUFFICIENCY-CHECK is polynomial in $N=|S|$.
\item Under ETH, there exist succinctly encoded instances with $k^* = \Omega(n)$ (indeed $k^*=n$) for which SUFFICIENCY-CHECK requires $2^{\Omega(n)}$ time.
\end{itemize}
For Boolean coordinate spaces ($N = 2^n$), the explicit-state bound is polynomial in $2^n$ (exponential in $n$), while under ETH the succinct lower bound yields $2^{\Omega(n)}$ time for the hard family in which $k^* = \Omega(n)$.
\end{corollary}

\begin{proposition}[Finite-State Query Lower-Bound Core via Empty-Set Subproblem]\label{prop:query-regime-obstruction}
In the finite-state query regime [Q\_fin], let $S$ be any finite state type with $|S|\ge 2$ and let $Q\subset S$ with $|Q|<|S|$.
Then there exist two decision problems $\mathcal{D}_{\mathrm{yes}},\mathcal{D}_{\mathrm{no}}$ over the same state space such that:
\begin{itemize}
\item their oracle views on all states in $Q$ are identical,
\item $\emptyset$ is sufficient for $\mathcal{D}_{\mathrm{yes}}$,
\item $\emptyset$ is not sufficient for $\mathcal{D}_{\mathrm{no}}$.
\end{itemize}
Consequently, no deterministic query procedure using fewer than $|S|$ state queries can solve SUFFICIENCY-CHECK on all such instances; the worst-case query complexity is $\Omega(|S|)$.
\leanmeta{\LH{CC38},
\LH{L6},
\LH{L12}.}
\end{proposition}

\begin{proof}
This is exactly the finite-state indistinguishability theorem
\LH{CC38}. For any $|Q|<|S|$, there is an unqueried hidden state $s_0$ producing oracle-indistinguishable yes/no instances with opposite truth values on the $I=\emptyset$ subproblem. Since SUFFICIENCY-CHECK contains that subproblem, fewer than $|S|$ queries cannot be correct on both instances, yielding the $\Omega(|S|)$ worst-case lower bound.
\end{proof}

\begin{corollary}[Information Barrier in Query Regimes]\label{cor:information-barrier-query}
For the fixed sufficiency question, finite query interfaces induce an information barrier:
\begin{itemize}
\item [Q\_fin] (Opt-oracle): with $|Q|<|S|$, there exist yes/no instances indistinguishable on all queried states but with opposite truth values on the $I=\emptyset$ subproblem.
\item [Q\_bool] value-entry and state-batch interfaces preserve the same obstruction at cardinality $<2^n$.
\end{itemize}
Hence the barrier is representational-access based: without querying enough of the hidden state support, exact certification is blocked by indistinguishability rather than by search-strategy choice.
\leanmeta{\LH{CC20},
\LH{CC22},
\LH{CC21}.}
\end{corollary}

\begin{corollary}[Boolean-Coordinate Instantiation]\label{cor:query-obstruction-bool}
In the Boolean-coordinate state space $S=\{0,1\}^n$, Proposition~\ref{prop:query-regime-obstruction} yields the familiar $\Omega(2^n)$ worst-case query lower bound for Opt-oracle access.
\leanmeta{\LH{CC37};
core: \LH{L5}.}
\end{corollary}

\begin{proposition}[Value-Entry Query Lower Bound]\label{prop:query-value-entry-lb}
In the mechanized Boolean value-entry query submodel [Q\_bool], for any deterministic procedure using fewer than $2^n$ value-entry queries $(a,s)\mapsto U(a,s)$, there exist two queried-value-indistinguishable instances with opposite truth values for SUFFICIENCY-CHECK on the $I=\emptyset$ subproblem. Therefore worst-case value-entry query complexity is also $\Omega(2^n)$.
\leanmeta{\LH{L8},
\LH{L17}.}
\end{proposition}

\begin{proof}
The theorem \LH{L8} constructs, for any query set of cardinality $<2^n$, an unqueried hidden state $s_0$ and a yes/no instance pair that agree on all queried values but disagree on $\emptyset$-sufficiency. The auxiliary bound \LH{L17} ensures that fewer than $2^n$ value-entry queries cannot cover all states, so the indistinguishability argument applies in the worst case.
\end{proof}

\begin{proposition}[Subproblem-to-Full Transfer Rule]\label{prop:query-subproblem-transfer}
If every full-problem solver induces a solver for a fixed subproblem, then any lower bound for that subproblem lifts to the full problem.
\leanmeta{\LH{CC61},
\LH{CC62},
\LH{CC42}.}
\end{proposition}

This is the restriction-map instance of Definition~\ref{def:regime-simulation}: a solver for the full regime induces one for the restricted subproblem regime, so lower bounds transfer.

\begin{proposition}[Randomized Robustness (Seedwise)]\label{prop:query-randomized-robustness}
In [Q\_bool], for any query set with cardinality $<2^n$, the indistinguishable yes/no pair from Proposition~\ref{prop:query-regime-obstruction} forces one decoding error \emph{per random seed} for any seed-indexed decoder from oracle transcripts.
Consequently, finite-support randomization does not remove the obstruction: averaging preserves a constant error floor on the hard pair.
\leanmeta{\LH{L10},
\LH{L11},
\LH{L4}.}
\end{proposition}

\begin{proposition}[Randomized Robustness (Weighted Form)]\label{prop:query-randomized-weighted}
For any finite-support seed weighting $\mu$, the same hard pair satisfies a weighted identity:
the weighted sum of yes-error and no-error equals total seed weight. Hence randomization cannot collapse both errors simultaneously.
\leanmeta{\LH{L21},
\LH{L22}.}
\end{proposition}

\begin{proposition}[State-Batch Query Lower Bound]\label{prop:query-state-batch-lb}
In [Q\_bool], the same $\Omega(2^n)$ lower bound holds for a state-batch oracle that returns the full Boolean-action utility tuple at each queried state.
\leanmeta{\LH{L7},
\LH{L13}.}
\end{proposition}

\begin{proposition}[Finite-State Empty-Subproblem Generalization]\label{prop:query-finite-state-generalization}
The empty-subproblem indistinguishability lower-bound core extends from Boolean-vector state spaces to any finite state type with at least two states.
\leanmeta{\LH{L6},
\LH{L12}.}
\end{proposition}

\begin{proposition}[Adversary-Family Tightness by Full Scan]\label{prop:query-tightness-full-scan}
For the const/spike adversary family used in the query lower bounds, querying all states distinguishes the pair; thus the lower-bound family is tight up to constant factors under full-state scan.
\leanmeta{\LH{L9}.}
\end{proposition}

\begin{proposition}[Weighted Query-Cost Transfer]\label{prop:query-weighted-transfer}
Let $w(q)$ be per-query cost and $w_{\min}$ a lower bound on all queried costs.
Any cardinality lower bound $|Q|\geq L$ lifts to weighted cost:
\[
\sum_{q\in Q} w(q)\ \ge\ w_{\min}\cdot L.
\]
\leanmeta{\LH{L19},
\LH{L20}.}
\end{proposition}

\begin{proposition}[Oracle-Lattice Transfer: Batch $\Rightarrow$ Entry]\label{prop:oracle-lattice-transfer}
In [Q\_bool], agreement on state-batch oracle views over touched states implies agreement on all value-entry views for the corresponding entry-query set.
\leanmeta{\LH{L18},
\LH{CC31}.}
\end{proposition}

This is the oracle-transducer instance of Definition~\ref{def:regime-simulation}: batch transcripts induce entry transcripts while preserving indistinguishability for the target sufficiency question.

\begin{proposition}[Oracle-Lattice Strictness: Opt vs Value Entries]\label{prop:oracle-lattice-strict}
`Opt`-oracle views are strictly coarser than value-entry views: there exist instances with identical `Opt` views but distinguishable value entries.
\leanmeta{\LH{L2},
\LH{L3}.}
\end{proposition}

\begin{remark}[Instantiation of Definitions~\ref{def:structural-complexity} and~\ref{def:representational-hardness}]
Theorem~\ref{thm:dichotomy} and Propositions~\ref{prop:query-regime-obstruction}--\ref{prop:query-value-entry-lb} keep the structural problem fixed (same sufficiency relation) and separate representational hardness by access regime: explicit-state access exposes the boundary $s \mapsto \Opt(s)$, finite-state query access already yields $\Omega(|S|)$ lower bounds at the Opt-oracle level (Boolean instantiation: $\Omega(2^n)$), value-entry/state-batch interfaces preserve the obstruction in [Q\_bool], and succinct access can hide structure enough to force ETH-level worst-case cost on a hard family.
\end{remark}

This regime-typed landscape identifies tractability under [E], a finite-state query lower-bound core under [Q\_fin] with Boolean interface refinements under [Q\_bool], and worst-case intractability under [S+ETH] for the same underlying decision question.

The structural reason for this separation is enumerable access. Under the explicit-state encoding, the mapping $s \mapsto \Opt(s)$ is directly inspectable and sufficiency verification reduces to scanning state pairs, at cost polynomial in $|S|$. Under succinct encoding, the circuit compresses an exponential state space into polynomial size; universal verification over that compressed domain---``for all $s, s'$ with $s_I = s'_I$, does $\Opt(s) = \Opt(s')$?''---carries worst-case cost exponential in $n$, because the domain the circuit compressed away cannot be reconstructed without undoing the compression.

\subsection{Budgeted Physical Crossover}
The [E] vs [S] split can be typed against an explicit physical budget. Let $E(n)$ and $S(n)$ denote explicit and succinct representation sizes for the same decision question at scale parameter $n$, and let $B$ be a hard representation budget.

\begin{proposition}[Budgeted Crossover Witness]\label{prop:budgeted-crossover}
If $E(n) > B$ and $S(n) \le B$ for some $n$, then explicit representation is infeasible at $(B,n)$ while succinct representation remains feasible at $(B,n)$.
\leanmeta{\LH{CC33},
\LH{PBC2}.}
\end{proposition}

\begin{proof}
This is exactly the definitional split encoded in \LH{PBC1}: explicit infeasibility is $B < E(n)$ and succinct feasibility is $S(n)\le B$. The theorem \LH{CC33} returns both conjuncts.
\end{proof}

\begin{proposition}[Above-Cap Crossover Existence]\label{prop:crossover-above-cap}
Assume there is a global succinct-size cap $C$ and explicit size is unbounded:
\[
\forall n,\ S(n)\le C,\qquad
\forall B,\ \exists n,\ B < E(n).
\]
Then for every budget $B\ge C$, there exists a crossover witness at budget $B$.
\leanmeta{\LH{CC32},
\LH{PBC3}.}
\end{proposition}

\begin{proof}
Fix $B\ge C$. By unboundedness of $E$, choose $n$ with $B < E(n)$. By the cap assumption, $S(n)\le C\le B$, so $(B,n)$ is a crossover witness. The mechanized closure theorem is \LH{CC32}.
\end{proof}

\begin{proposition}[Crossover Does Not Imply Exact-Certification Competence]\label{prop:crossover-not-certification}
Fix any collapse schema ``exact certification competence $\Rightarrow$ complexity collapse'' for the target exact-certification predicate. Under the corresponding non-collapse assumption, exact-certification competence is impossible even when a budgeted crossover witness exists.
\leanmeta{\LH{CC34}.}
\end{proposition}

\begin{proof}
The mechanized theorem \LH{CC34} bundles the crossover split with the same logical closure used in Proposition~\ref{prop:integrity-resource-bound}: representational feasibility and certification hardness are distinct predicates.
\end{proof}

\begin{proposition}[Crossover Policy Closure Under Integrity]\label{prop:crossover-policy}
In the certifying-solver model (Definitions~\ref{def:solver-integrity}--\ref{def:competence-regime}), assume:
\begin{itemize}
\item a budgeted crossover witness at $(B,n)$,
\item the same collapse implication as Proposition~\ref{prop:integrity-resource-bound},
\item a solver that is integrity-preserving for the declared relation.
\end{itemize}
Then either full in-scope non-abstaining coverage fails, or declared budget compliance fails.
Equivalently, in uncertified hardness-blocked regions, integrity is compatible with $\mathsf{ABSTAIN}/\mathsf{UNKNOWN}$ but not with unconditional exact-competence claims.
\leanmeta{\LH{CC35}.}
\end{proposition}

\subsection{Thermodynamic Lift in Declared Model}
The crossover and hardness statements are complexity-theoretic.
The thermodynamic lift follows the same assumption discipline as the rest of the paper:
claims are exactly those derivable from the declared conversion model and declared bit-operation lower-bound premises.

\begin{proposition}[Bit-to-Energy/Carbon Lift]\label{prop:thermo-lift}
Fix a declared thermodynamic conversion model with constants $(\lambda,\rho)$ where $\lambda$ lower-bounds energy per irreversible bit operation and $\rho$ lower-bounds carbon per joule. If a bit-operation lower bound $b_{\mathrm{LB}} \le b_{\mathrm{used}}$ holds, then:
\[
E_{\mathrm{LB}}=\lambda b_{\mathrm{LB}} \le \lambda b_{\mathrm{used}},\qquad
C_{\mathrm{LB}}=\rho E_{\mathrm{LB}} \le \rho(\lambda b_{\mathrm{used}}).
\]
\leanmeta{\LH{CC73},
\LH{CC74}.}
\end{proposition}

\begin{proposition}[Hardness--Thermodynamics Bundle]\label{prop:thermo-hardness-bundle}
Under the same non-collapse/collapse schema used by Proposition~\ref{prop:integrity-resource-bound}, exact-certification competence is impossible; combined with a declared bit-operation lower bound, this yields simultaneous lower bounds on energy and carbon in the declared conversion model.
\leanmeta{\LH{CC75}.}
\end{proposition}

\begin{proposition}[Mandatory Physical Cost]\label{prop:thermo-mandatory-cost}
Within the declared linear thermodynamic model, if
\[
\lambda > 0,\quad \rho > 0,\quad b_{\mathrm{LB}} > 0,
\]
then both lower bounds are strictly positive:
\[
E_{\mathrm{LB}}=\lambda b_{\mathrm{LB}} > 0,\qquad
C_{\mathrm{LB}}=\rho E_{\mathrm{LB}} > 0.
\]
\leanmeta{\LH{CC76}.}
\end{proposition}

\begin{proposition}[Conserved Additive Accounting]\label{prop:thermo-conservation-additive}
Within the same declared linear model, lower-bound accounting is additive over composed bit-operation lower bounds:
\[
E_{\mathrm{LB}}(b_1+b_2)=E_{\mathrm{LB}}(b_1)+E_{\mathrm{LB}}(b_2),\qquad
C_{\mathrm{LB}}(b_1+b_2)=C_{\mathrm{LB}}(b_1)+C_{\mathrm{LB}}(b_2).
\]
\leanmeta{\LH{CC72}.}
\end{proposition}
